\begin{frame}{프로젝트 구조 — 환경 · 적용 · 필수 요구사항}
  \begin{itemize}
    \item \textbf{환경}: PyBullet 또는 Gymnasium의 연속제어 환경
      (Pendulum, MuJoCo 기반 로봇 등) 중 하나.
    \item \textbf{적용}: Week 9--15 중 \textbf{최소 하나의 방법}(DQN, PPO,
      모방학습, MCTS)으로 정책을 학습.
    \item \textbf{필수 — 수식적 정당화 최소 1개}, 예:
      \begin{itemize}
        \item PPO: Ch11.3의 클리핑이 학습 안정성에 미친 영향(있음/없음 비교)
        \item 모방학습: Ch12.1의 복합 오차 관찰(전문가 데이터 분포 밖 성능 저하)
        \item 도메인 무작위화: Ch13.1의 현실 격차 논의와 연결해 효과 설명
      \end{itemize}
    \item \kb{학습 곡선(에피소드별 누적 보상)은 반드시 제시} — 지도학습의
      손실 곡선에 해당하는, ``학습이 실제로 되고 있는가''의 기본 증거.
  \end{itemize}
\end{frame}
